% Chapter 4, Topic _Linear Algebra_ Jim Hefferon
%  http://joshua.smcvt.edu/linearalgebra
%  2001-Jun-12
\topic{幂方法}
\index{method of powers|(}
\index{powers, method of|(}
在应用中，矩阵可能非常庞大。
通过求出并求解特征多项式来计算特征值与特征向量并不现实，
既太慢又容易出错。
有些技巧可以避开特征多项式。
本节将介绍一种适用于大型矩阵的方法，这样的矩阵是
\definend{稀疏}\index{matrix!sparse}\index{sparse matrix} 的，
也就是说，其中绝大多数元素都是零。

设 $\nbyn{n}$ 矩阵 $T$ 有 $n$ 个互不相同的特征值
$\lambda_1$、$\lambda_2$、\ldots、$\lambda_n$。
那么 $\C^n$ 有一组由相应的特征向量
$\sequence{\vec{\zeta}_1,\dots,\vec{\zeta}_n}$ 构成的基。
对任意 $\vec{v}\in\C^n$，记
$\vec{v}=c_1\vec{\zeta}_1+\dots+c_n\vec{\zeta}_n$，对 $\vec{v}$ 反复作用以 $T$
便得到下面各式。
\begin{align*}
  T\vec{v}
      &=c_1\lambda_1\vec{\zeta}_1+c_2\lambda_2\vec{\zeta}_2+
                              \dots+c_n\lambda_n\vec{\zeta}_n  \\
  T^2\vec{v}
      &=c_1\lambda_1^2\vec{\zeta}_1+c_2\lambda_2^2\vec{\zeta}_2+
                              \dots+c_n\lambda_n^2\vec{\zeta}_n  \\
  T^3\vec{v}
      &=c_1\lambda_1^3\vec{\zeta}_1+c_2\lambda_2^3\vec{\zeta}_2+
                              \dots+c_n\lambda_n^3\vec{\zeta}_n  \\
      &\vdotswithin{=}                                            \\
  T^k\vec{v}
      &=c_1\lambda_1^k\vec{\zeta}_1+c_2\lambda_2^k\vec{\zeta}_2+
                              \dots+c_n\lambda_n^k\vec{\zeta}_n
\end{align*}
假设 $|\lambda_1|$
是其中最大的，将两边同除以 $\lambda_1^k$
\begin{equation*}
  \frac{T^k\vec{v}}{\lambda_1^k}
      =c_1\vec{\zeta}_1+c_2\frac{\lambda_2^k}{\lambda_1^k}\vec{\zeta}_2+
                    \dots+c_n\frac{\lambda_n^k}{\lambda_1^k}\vec{\zeta}_n
\end{equation*}
可知，随着 $k$ 的增大，各分式趋于零，
于是 $\lambda_1$ 的项将在表达式中占主导地位，
而该表达式的极限为 $c_1\vec{\zeta}_1$。

于是若 $c_1\neq 0$，
则随着 $k$ 增大，向量 $T^k\vec{v}$ 将趋于与主特征值相关联的特征向量的方向。
从而，
向量长度之比 $\absval{T^{k}\vec{v}}/\absval{T^{k-1}\vec{v}}$
趋于该主特征值。

例如，矩阵
\begin{equation*}
  T=\begin{mat}[r]
    3  &0  \\
    8  &-1
  \end{mat}
\end{equation*}
的特征值是 $3$ 和 $-1$。
若 $\vec{v}$ 的分量为 $1$ 和 $1$，则迭代给出如下结果。
\begin{center}
  \begin{tabular}{c|ccccc}
     $\vec{v}$  &$T\vec{v}$  &$T^2\vec{v}$
        &$\cdots$ &$T^9\vec{v}$ &$T^{10}\vec{v}$        \\ \hline
     $\matrixvenlarge{\colvec[r]{1 \\ 1}}$
        &$\matrixvenlarge{\colvec[r]{3 \\ 7}}$
        &$\matrixvenlarge{\colvec[r]{9 \\ 17}}$
        &$\cdots$
        &$\matrixvenlarge{\colvec[r]{19\,683 \\ 39\,367}}$
        &$\matrixvenlarge{\colvec[r]{59\,049 \\ 118\,097}}$
  \end{tabular}
\end{center}
最后两个的长度之比是 $2.999\,9$。

我们指出两个实现上的问题。
其一，
我们不先求出 $T$ 的各次幂再作用到 $\vec{v}$ 上，
而是先把 $\vec{v}_1$ 算成 $T\vec{v}$，再把 $\vec{v}_2$ 算成
$T\vec{v}_1$，依此类推（也就是说，我们并不单独计算
$T^2\!$、$T^3\!$、\ldots）。
即使 $T$ 很大，只要它是稀疏的，
这些矩阵与向量的乘积也能很快完成。
第二个问题是，
为避免产生的数大到超出计算机的能力而溢出，
可以在每一步对各个 $\vec{v}_i$ 作归一化。
例如，可以把每个 $\vec{v}_i$ 除以它的长度
（其他做法还有：除以它的最大分量，或干脆除以它的第一个分量）。
于是我们通过生成下述序列来实现这一方法
\begin{align*}
  \vec{w}_0  &=\vec{v}_0/\absval{\vec{v}_0} \\
  \vec{v}_1  &=T\vec{w}_0                 \\
  \vec{w}_1  &=\vec{v}_1/\absval{\vec{v}_1} \\
  \vec{v}_2  &=T\vec{w}_2                 \\
             &\vdotswithin{=}    \\
  \vec{w}_{k-1}  &=\vec{v}_{k-1}/\absval{\vec{v}_{k-1}} \\
  \vec{v}_k  &=T\vec{w}_k
\end{align*}
直到满意为止。
此时 $\vec{v}_k$ 是某个特征向量的一个近似，而
主特征值的近似是
比值 % $\absval{\vec{v}_{k}}/\absval{\vec{w}_{k-1}}$.
$(T\dotprod \vec{v}_k)/(\vec{v}_k\dotprod\vec{v}_k)
  \approx (\lambda_1\vec{v}_k\cdot\vec{v}_k)/(\vec{v}_k\dotprod\vec{v}_k)
  =\lambda_1$。


判断“满意”的一种方式是迭代到特征值的近似值稳定下来为止。
例如可以决定，迭代的终止不设在某个固定的步数之后，
而是设在 $\absval{\vec{v}_k}$ 与 $\absval{\vec{v}_{k-1}}$
相差小于百分之一时，或者它们到第二位有效数字都一致时。

收敛的速度由 $|\lambda_2/\lambda_1|$ 的幂趋于零的速度决定，
其中 $\lambda_2$ 是长度第二大的特征值。
若该比值远小于一，则收敛很快；
但若它只比一小一点，收敛就可能相当慢。
因此，幂方法
并不是求特征值最常用的方法
（尽管它是最简单的一种，这正是它出现在这里的原因）。
作为替代，有一大类各种各样的方法，它们通常这样工作：先把给定的矩阵 $T$
换成另一个与它相似
从而有相同特征值的矩阵，但它处于某种简化形式，
例如
三对角形式\index{tridiagonal form}\index{matrix!tridiagonal}，
即除对角线上以及紧贴对角线上、下两条斜线上的元素外，
其余元素全为零。
然后可以用针对特殊形式的技巧求出特征值。
一旦知道特征值，便能容易地算出
$T$ 的特征向量。
这些其他方法超出了本书的范围。
一本好的参考书是 \cite{Goult}



\begin{exercises}
  \item
    用十次迭代估计下列矩阵的最大特征值，
    初始向量取分量为 $1$ 和 $2$ 的向量。
    将所得答案与通过求解特征方程
    得到的答案比较。
    \begin{exparts*}
      \partsitem $\begin{mat}[r]
                    1  &5  \\
                    0  &4
                  \end{mat}$
      \partsitem $\begin{mat}[r]
                    3   &2  \\
                    -1  &0
                  \end{mat}$
    \end{exparts*}
    \begin{answer}
     \begin{exparts}
       \partsitem 一眼即可看出，最大特征值是 $4$。
         \Sage{} 给出如下结果。
\begin{lstlisting}
sage: def eigen(M,v,num_loops=10):
....:     for p in range(num_loops):
....:         v_normalized = (1/v.norm())*v
....:         v = M*v
....:     return v
....:
sage: M = matrix(RDF, [[1,5], [0,4]])
sage: v = vector(RDF, [1, 2])
sage: v = eigen(M,v)
sage: (M*v).dot_product(v)/v.dot_product(v)
4.00000147259
\end{lstlisting}
       \partsitem 简单的计算表明，最大特征值
          为~$2$。
          \Sage{} 给出如下结果。
\begin{lstlisting}
sage: M = matrix(RDF, [[3,2], [-1,0]])
sage: v = vector(RDF, [1, 2])
sage: v = eigen(M,v)
sage: (M*v).dot_product(v)/v.dot_product(v)
2.00097741083
\end{lstlisting}
     \end{exparts}
    \end{answer}
  \item
     重做前面的练习：迭代直到
     $\absval{\vec{v}_k}-\absval{\vec{v}_{k-1}}$ 的绝对值小于
     $0.01$。
     每一步都将每个向量除以其长度来归一化。
     这需要多少次迭代？
     两种答案的差异显著吗？
     \begin{answer}
       \begin{exparts}
         \partsitem 下面是 \Sage{} 的结果。
 \begin{lstlisting}
sage: def eigen_by_iter(M, v, toler=0.01):
....:     dex = 0
....:     diff = 10
....:     while abs(diff)>toler:
....:         dex = dex+1
....:         v_next = M*v
....:         v_normalized = (1/v.norm())*v
....:         v_next_normalized = (1/v_next.norm())*v_next
....:         diff = (v_next_normalized-v_normalized).norm()
....:         v_prior = v_normalized
....:         v = v_next_normalized
....:     return v, v_prior, dex
....:
sage: M = matrix(RDF, [[1,5], [0,4]])
sage: v = vector(RDF, [1, 2])
sage: v,v_prior,dex = eigen_by_iter(M,v)
sage: (M*v).norm()/v.norm()
4.00604111686
sage: dex
4
 \end{lstlisting}
      \partsitem 对这道题 \Sage{} 要多迭代几次。
        这里用到了上一小题中定义的过程。
\begin{lstlisting}
sage: M = matrix(RDF, [[3,2], [-1,0]])
sage: v = vector(RDF, [1, 2])
sage: v,v_prior,dex = eigen_by_iter(M,v)
sage: (M*v).norm()/v.norm()
2.01585174302
sage: dex
6
\end{lstlisting}
       \end{exparts}
     \end{answer}
  \item
    用十次迭代估计下列矩阵的最大特征值，
    初始向量取分量为 $1$、$2$、$3$ 的向量。
    将所得答案与通过求解特征方程
    得到的答案比较。
    \begin{exparts*}
      \partsitem $\begin{mat}[r]
                    4   &0  &1 \\
                    -2  &1  &0  \\
                    -2  &0  &1
                  \end{mat}$
      \partsitem $\begin{mat}[r]
                   -1  &2  &2  \\
                    2  &2  &2  \\
                   -3  &-6 &-6
                  \end{mat}$
    \end{exparts*}
    \begin{answer}
     \begin{exparts}
       \partsitem 最大特征值是 $3$。
         \Sage{} 给出如下结果。
\begin{lstlisting}
sage: M = matrix(RDF, [[4,0,1], [-2,1,0], [-2,0,1]])
sage: v = vector(RDF, [1, 2, 3])
sage: v = eigen(M,v)
sage: (M*v).dot_product(v)/v.dot_product(v)
3.02362112326
\end{lstlisting}
       \partsitem 最大特征值是 $-3$。
\begin{lstlisting}
sage: M = matrix(RDF, [[-1,2,2], [2,2,2], [-3,-6,-6]])
sage: v = vector(RDF, [1, 2, 3])
sage: v = eigen(M,v)
sage: (M*v).dot_product(v)/v.dot_product(v)
-3.00941127145
\end{lstlisting}
     \end{exparts}
    \end{answer}
  \item
     重做前面的练习：迭代直到
     $\absval{\vec{v}_k}-\absval{\vec{v}_{k-1}}$ 的绝对值小于
     $0.01$。
     每一步都将每个向量除以其长度来归一化。
     这需要多少次迭代？
     两种答案的差异显著吗？
     \begin{answer}
       \begin{exparts}
         \partsitem
           \Sage{} 得到如下结果，其中 \lstinline!eigen_by_iter!
           已在上面定义。
\begin{lstlisting}
sage: M = matrix(RDF, [[4,0,1], [-2,1,0], [-2,0,1]])
sage: v = vector(RDF, [1, 2, 3])
sage: v,v_prior,dex = eigen_by_iter(M,v)
sage: (M*v).dot_product(v)/v.dot_product(v)
3.05460392934
sage: dex
8
\end{lstlisting}
         \partsitem 采用如下设置，
\begin{lstlisting}
sage: M = matrix(RDF, [[-1,2,2], [2,2,2], [-3,-6,-6]])
sage: v = vector(RDF, [1, 2, 3])
\end{lstlisting}
            \Sage{} 不返回结果（可用 \lstinline!<Ctrl>-c! 中断
            计算）。
            在该例程中加入一些错误检查代码
\begin{lstlisting}
def eigen_by_iter(M, v, toler=0.01):
    dex = 0
    diff = 10
    while abs(diff)>toler:
        dex = dex+1
        if dex>1000:
            print "oops! probably in some loop: \nv=",v,"\nv_next=",v_next
        v_next = M*v
        if (v.norm()==0):
            print "oops! v is zero"
	    return None
        if (v_next.norm()==0):
            print "oops! v_next is zero"
            return None
        v_normalized = (1/v.norm())*v
        v_next_normalized = (1/v_next.norm())*v_next
        diff = (v_next_normalized-v_normalized).norm()
        v_prior = v_normalized
        v = v_next_normalized
    return v, v_prior, dex
\end{lstlisting}
           得到如下输出。
\begin{lstlisting}
oops! probably in some loop:
  v= (0.707106781187, -1.48029736617e-16, -0.707106781187)
  v_next= (2.12132034356, -4.4408920985e-16, -2.12132034356)
oops! probably in some loop:
  v= (-0.707106781187, 1.48029736617e-16, 0.707106781187)
  v_next= (-2.12132034356, 4.4408920985e-16, 2.12132034356)
oops! probably in some loop:
  v= (0.707106781187, -1.48029736617e-16, -0.707106781187)
  v_next= (2.12132034356, -4.4408920985e-16, -2.12132034356)
\end{lstlisting}
          可见它在循环打转。
       \end{exparts}
     \end{answer}
  \item
      若 $c_1=0$ 会怎样？
      也就是说，若初始向量在相应特征向量的方向上
      没有任何分量，会出现什么情况？
     \begin{answer}
       理论上，这一方法将给出 $\lambda_2$。
       然而在实际计算中，舍入误差会引入
       $\vec{v}_1$ 方向上的分量，因此该方法
       仍会给出 $\lambda_1$，只是与选取较为幸运的初始向量
       相比，可能要多花些时间。
     \end{answer}
  \item
    如何把幂方法加以改造，
    用来求最小的特征值？
    \begin{answer}
      不使用 $\vec{v}_k=T\vec{v}_{k-1}$，
      而使用 $T^{-1}\vec{v}_k=\vec{v}_{k-1}$。
    \end{answer}
\end{exercises}

\announcecomputercode
下面是完成上述计算的计算机代数系统 Octave 的代码。
（已稍作编辑，删去了空行等。）
\begin{lstlisting}
>T=[3, 0;
    8, -1]
T=
   3   0
   8  -1
>v0=[1; 2]
v0=
   1
   1
>v1=T*v0
v1=
   3
   7
>v2=T*v1
v2=
   9
  17
>T9=T**9
T9=
  19683  0
  39368 -1
>T10=T**10
T10=
  59049  0
 118096  1
>v9=T9*v0
v9=
  19683
  39367
>v10=T10*v0
v10=
  59049
 118096
>norm(v10)/norm(v9)
ans=2.9999
\end{lstlisting}
\textit{注。}
这里没有用到 Octave 的全部本领；
它内置了若干函数，可自动运用复杂的方法
求出特征值与特征向量。

\index{powers, method of|)}
\index{method of powers|)}
\endinput





